\documentclass{article}

% ============================================================
%  PACKAGE OPTIONS
% ============================================================
% \usepackage{neurips_2026}              % submission (anonymous)
% \usepackage[preprint]{neurips_2026}    % preprint
\usepackage[final]{neurips_2026}         % camera-ready / final

% ---------------- standard packages ----------------
\usepackage[utf8]{inputenc}
\usepackage[T1]{fontenc}
\usepackage{hyperref}
\usepackage{url}
\usepackage{booktabs}
\usepackage{amsfonts}
\usepackage{amsmath}
\usepackage{nicefrac}
\usepackage{microtype}
\usepackage{xcolor}
\usepackage{multirow}
\usepackage{graphicx}
\usepackage{array}
\usepackage{subcaption}
\usepackage{siunitx}
\usepackage{placeins}

\captionsetup{font=footnotesize, skip=3pt}
\setlength{\textfloatsep}{10pt plus 2pt minus 2pt}
\setlength{\floatsep}{8pt plus 2pt minus 2pt}
\setlength{\intextsep}{8pt plus 2pt minus 2pt}

% ---------------- float placement (loosened to avoid big gaps) ----------------
\renewcommand{\topfraction}{0.9}
\renewcommand{\bottomfraction}{0.8}
\renewcommand{\textfraction}{0.1}
\renewcommand{\floatpagefraction}{0.8}
\setcounter{topnumber}{3}
\setcounter{bottomnumber}{3}
\setcounter{totalnumber}{5}

\makeatletter
\def\@trackname{}
\makeatother

% ---------------- title / author ----------------
\title{Your Report Title Here}

\author{
    Muhammad Aitasaam Rana \\
    Roll Number: 28100374 \\
    \texttt{28100374@lums.edu.pk} \\
}

\begin{document}

\maketitle

\begin{abstract}
Write a short summary of the problem, what you did, and your main findings (150--250 words).

Code for all experiments is available at:
\url{https://github.com/your-username/your-repo}.
\end{abstract}

% ============================================================
\section{Introduction}
% ============================================================
State the core question, motivation, and a brief overview of the tasks or experiments. Cite related work with \citep{key} or \citet{key}.

% ============================================================
\section{Experimental Setup / Methods}
% ============================================================
% ------------------------------------------------------------
\subsection{Task 1}
\label{sec:task1-protocol}
% ------------------------------------------------------------

\paragraph{Setup and models.}
We fine-tune \texttt{Qwen3-0.6B} on $40{,}000$ MetaMathQA examples (an LLM-expanded version of GSM8K and MATH problems) and evaluate seven checkpoints: (1)~the untouched \textbf{Baseline}; (2)~\textbf{Full FT Midpoint} ($625$ updates, $20{,}000$ examples); (3)~\textbf{Full FT Full} ($1{,}250$ updates, 1 epoch); (4)~\textbf{Full FT + $L_2$-SP}, full fine-tuning with $L_2$-SP distance regularization; (5)~\textbf{LoRA-4} and (6)~\textbf{LoRA-16}, standard LoRA with rank $4$ and $16$; and (7)~\textbf{OPLoRA}, double-sided Orthogonal Projection LoRA (final checkpoint). Training sequences are capped at $2{,}048$ tokens without truncating target answers. To keep evaluation strictly zero-shot, any prompt shared by MetaMathQA and the GSM8K test set (after text normalization) is removed from training.

\paragraph{Evaluation.}
All models are evaluated zero-shot. \textbf{GSM8K} (math transfer) uses generation: the final numerical answer is parsed by regex, and unextractable outputs count as incorrect. Generating also lets us inspect response style and whether reasoning changes; the cap is $1{,}024$ new tokens, half the training cap. \textbf{HellaSwag} (commonsense) and \textbf{ARC-Easy} (science) use candidate likelihood, reporting raw (\texttt{acc}) and length-normalized (\texttt{acc\_norm}) accuracy. \textbf{BoolQ} (binary reading) also uses candidate likelihood; since yes/no options have equal token length, only raw accuracy is reported.

\paragraph{Diagnostics.}
\textbf{Activation probes} compare each method's hidden-layer activations against the baseline on $64$ fixed prompts ($32$ GSM8K test and $32$ HellaSwag validation samples, seed $6304$). A \textbf{qualitative non-math inspection} uses $24$ fixed prompts from \texttt{first\_plot\_questions.yaml} ($8$ core questions in plain, JSON, and template formats), reporting instruction compliance, coherence, answer style, and inappropriate content.

\paragraph{Checkpointing.}
To stay under Kaggle's $20$\,GB disk limit, we kept only the most recent recovery checkpoint on disk, overwriting older ones as training progressed. For the final step-$1{,}250$ checkpoint we saved only the model weights and dropped the optimizer state, since it is needed only to resume training, not for benchmark evaluation or diagnostic probes.

\paragraph{Task 2: Name.}
Method description.

% ============================================================
\section{Results}
% ============================================================

% ------------------------------------------------------------
\subsection{Task 1: Name}
\label{sec:task1-results}
% ------------------------------------------------------------

Discussion of results, referring to Table~\ref{tab:example} and Figure~\ref{fig:example}.


% ------------------------------------------------------------
\subsection{Task 2: Name}
\label{sec:task2-results}
% ------------------------------------------------------------

Discussion of results.

\FloatBarrier

% ============================================================
\section{Cross-task Discussion}
% ============================================================

\subsection{Theme 1}
Synthesize findings across tasks.

\subsection{Theme 2}
Further synthesis.

% ============================================================
\section{Conclusion}
% ============================================================
Summarize findings, state limitations (training budget, dataset size, seeds), and suggest future work.

% ============================================================
%  REFERENCES
% ============================================================
% \nocite{*}   % uncomment to list every entry in references.bib
\bibliographystyle{plainnat}
\bibliography{references}

% ============================================================
%  APPENDIX
% ============================================================
\newpage
\appendix

\section{Task 1: Supplementary Material}
\label{app:task1}

Additional tables and figures.

\FloatBarrier

\end{document}